
This appendix records how AlphaFold3 was brought up on each backend,
its timings, and the per-sample score comparisons summarized in
Section~\ref{sec:alphafold3}.

\subsection{Bringing AlphaFold3 up}
\label{app:af3-infra}

AlphaFold3 is a separate, diffusion-based codebase, not a version
increment of AlphaFold2, and required meaningfully more engineering to
run on every backend. Unlike AlphaFold2's pure-Python install,
AlphaFold3 needs a native C++ build step (\texttt{libcifpp} and pybind11
bindings) unavailable without root on our cluster's login node, and a
separate \texttt{build\_data} step that constructs the Chemical Component
Dictionary the model needs at start-up. Skipping it does not fail at
start-up but a few seconds into the run, with a missing-file error. The T4
also required \texttt{XLA\_FLAGS=-{}-xla\_disable\_hlo\_passes=%
custom-kernel-fusion-rewriter}, documented in AlphaFold3's own
Dockerfile comments: without it the run raises a \texttt{ValueError}
before any computation starts. Its
released weights decompress to 1.15\,GB. That is roughly 3.3$\times$ an
AlphaFold2 model, measured against an approximately 350\,MB AlphaFold2
figure taken from our own notes; we did not weigh a checkpoint, since
these runs initialize AlphaFold2 randomly and never load a trained
one.
On our cluster, the C++ build step had to run inside a Kubernetes Job,
for the same reason the rest of this paper's
TPU work does: no root access outside a Job.

\subsection{CPU and GPU results}
\label{sec:af3-cpuresults}

AlphaFold3's own timed model-inference region (one seed, five diffusion
samples, including JIT compilation, since AlphaFold3 makes a single
model call per process with no warm-up) runs 2401.38\,s on the Colab
CPU and 86.21\,s on the Colab T4, a 27.86$\times$ GPU-over-CPU speedup
with compilation included on both sides.
As with AlphaFold2, we do not compute an AlphaFold3-over-AlphaFold2
speed ratio anywhere in this paper: the two models differ in recycling
depth, in whether the timed region includes compilation, and in how
many output samples one call produces, and an early draft of this
comparison that did not control for those differences reached the
wrong sign. We report each model's own internal ratios and nothing that
divides one model's timing by the other's.

\subsection{A reproducibility finding}
\label{app:af3-repro}

Section~\ref{sec:af3-body-repro} gives the headline comparison. Across
backends, the best-ranked sample's score falls from 0.41 to 0.33 and its
fraction of disordered residues from 0.37 to 0.21.
The divergence is not explained by anything in our own methodology: the
two runs share seed, input and AlphaFold3 version. They are not identical in every respect: the GPU
run also sets the XLA flag described above, which the CPU run does not
need, and we did not isolate that flag's effect on the outputs.

A plausible, independently documented cause exists. AlphaFold3's own
performance documentation states that CUDA Capability 7.x devices have
known numerical issues, and recommends a specific XLA flag as a
workaround \citep{af3docs2026}; the Colab T4 used throughout this paper
is compute capability 7.5, inside that range, and the same flag was
required to get any AlphaFold3 result on it at all
(Appendix~\ref{app:af3-infra}). Community reports on AlphaFold3's own
issue tracker corroborate output problems on GPUs below the officially
tested tier (A100 and H100), though we did not find that tracker
discussion naming compute capability 7.5 specifically. We treat the
documented CUDA-capability issue, not the tracker discussion, as the
source for this explanation, and report it as a plausible cause
consistent with our own measurements. It is not confirmed: we did not
instrument AlphaFold3 internally to isolate which operation
produces the divergence.

